\section{Substructures}
\label{sec:maps-elementarity-02}

Substructure, elementary substructure, and the Tarski--Vaught test were
introduced for the first-order metatheory in
\cref{sec:fol-cardinality-normalforms-01}.  This chapter restates them once
in its map-based framework; the definitions agree, and the restatement is
recorded here rather than at each item.

\begin{definition}[Substructure]
The structure \(\mathcal M\) is a \emph{substructure} of \(\mathcal N\),
written \(\mathcal M\subseteq\mathcal N\), iff \(M\subseteq N\), the
inclusion map is an embedding, and therefore:
\begin{enumerate}
  \item every constant of \(\mathcal N\) lies in \(M\);
  \item \(M\) is closed under every function of \(\mathcal N\); and
  \item every relation of \(\mathcal M\) is the restriction to \(M\) of the
        corresponding relation of \(\mathcal N\).
\end{enumerate}
\end{definition}

Suppose that \(A\subseteq N\) is nonempty, or that \(\Lang\) has a constant
symbol.  The \emph{substructure generated by \(A\)} is the intersection of all
substructures of \(\mathcal N\) whose domains contain \(A\).  Its domain
consists exactly of values \(t^{\mathcal N}(\bar a)\), where \(t\) is a term
and \(\bar a\) is a finite tuple from \(A\), including the values of constant
terms.  If structures are required to be nonempty, \(A=\varnothing\), and the
signature has no constant, a least substructure need not exist; the phrase
``generated by \(\varnothing\)'' then requires a separate convention.

\begin{example}[A substructure that is not elementary]
In the relational signature \(\Lang=\{<\}\), the standard order
\((\N,<)\) is a substructure of \((\Z,<)\).  It is not elementary.  With the
parameter \(0\in\N\),
\[
  (\Z,<)\entails\exists y\,(y<0),
  \qquad
  (\N,<)\notentails\exists y\,(y<0).
\]
This example also shows why agreement on atomic formulas does not imply
agreement on quantified formulas.
\end{example}

\begin{lemma}[Isomorphism invariance of satisfaction]
\label{lem:iso-invariance}
Let \(h:\mathcal M\to\mathcal N\) be an isomorphism.  Then, for every
formula \(\varphi(x_1,\ldots,x_k)\) and all \(a_1,\ldots,a_k\in M\),
\[
  \mathcal M\entails\varphi(\bar a)
  \quad\Longleftrightarrow\quad
  \mathcal N\entails\varphi(h\bar a).
\]
In particular, isomorphic structures satisfy the same sentences.
\textit{Status: complete proof.}
\end{lemma}
\begin{proof}
Induct on \(\varphi\).  For an equality atom \(t_1=t_2\),
\cref{lem:term-preservation} gives
\(h(t_i^{\mathcal M}(\bar a))=t_i^{\mathcal N}(h\bar a)\) for both terms,
and injectivity of \(h\) yields
\(t_1^{\mathcal M}(\bar a)=t_2^{\mathcal M}(\bar a)\) iff
\(t_1^{\mathcal N}(h\bar a)=t_2^{\mathcal N}(h\bar a)\).  For a relation
atom \(R(t_1,\ldots,t_l)\), \cref{lem:term-preservation} reduces the claim
to the preservation and reflection of \(R\) required of an embedding.
The Boolean cases follow from the induction hypotheses and the
satisfaction clauses.  For \(\exists y\,\psi\): a witness \(a'\in M\) on
the left maps to the witness \(h(a')\) on the right by the induction
hypothesis; conversely, a witness \(b\in N\) equals \(h(a')\) for some
\(a'\in M\) by surjectivity, and the induction hypothesis returns the
witness \(a'\).  The universal case is dual, using the same surjectivity.
\end{proof}

\begin{lemma}[Image of an embedding]
\label{lem:embedding-image}
Let \(h:\mathcal M\to\mathcal N\) be an embedding.  Then \(h(M)\)
contains \(c^{\mathcal N}\) for every constant symbol \(c\) and is closed
under \(f^{\mathcal N}\) for every function symbol \(f\).  The structure
\(h[\mathcal M]\) with domain \(h(M)\), with the constants
\(c^{\mathcal N}\), with the restrictions of the functions
\(f^{\mathcal N}\), and with the relations \(R^{\mathcal N}\) restricted
to \(h(M)\) is a substructure of \(\mathcal N\), and \(h\), taken with
codomain \(h(M)\), is an isomorphism from \(\mathcal M\) to
\(h[\mathcal M]\).
\textit{Status: complete proof.}
\end{lemma}
\begin{proof}
By the homomorphism clauses,
\(c^{\mathcal N}=h(c^{\mathcal M})\in h(M)\), and
\(f^{\mathcal N}(h\bar a)=h(f^{\mathcal M}(\bar a))\in h(M)\) for every
tuple \(h\bar a\) from \(h(M)\); so \(h(M)\) is closed as claimed and the
definition of \(h[\mathcal M]\) is legitimate.  Its interpretations are by
construction the restrictions of those of \(\mathcal N\), so the
inclusion map preserves and reflects atomic formulas, and
\(h[\mathcal M]\subseteq\mathcal N\) in the sense of the substructure
definition above.  The map \(h\) with codomain \(h(M)\) is injective,
because \(h\) is, and surjective by construction.  It is a homomorphism
into \(h[\mathcal M]\), since the interpretations there are the
restrictions of those of \(\mathcal N\) to a set containing every
required value, and it reflects every relation symbol \(R\) because the
relation of \(h[\mathcal M]\) interpreting \(R\) is the restriction of
\(R^{\mathcal N}\) to tuples from \(h(M)\) while \(h\) reflects
\(R^{\mathcal N}\).  Hence the corestricted map is a surjective
embedding, which is an isomorphism by definition.
\end{proof}

